% \section{Objectives}
% \section{Challenges}
% \section{Contributions}

\chapter{Introduction}
\label{chap:introduction}

\section{Motivation}

Large Language Models (LLMs) have moved from being standalone text generators to being embedded inside increasingly complex software systems. Early progress in large-scale language modelling showed that Transformer-based models could perform a wide range of tasks from natural language prompts alone \cite{vaswani2017attention,brown2020language}, and instruction tuning made these models more useful as general-purpose assistants that follow user intent \cite{ouyang2022training}. More recently, LLMs have been used as agents: systems that plan over multiple steps, call tools, retrieve information, interact with users, and modify external state. Approaches such as ReAct demonstrated the value of combining reasoning with action \cite{yao2023react}, while work such as Toolformer and API-Bank showed the growing importance of external tool use in language model systems \cite{schick2023toolformer,li2023apibank}. This shift has changed the central reliability question. It is no longer sufficient to ask whether a model can produce a plausible final answer; we must ask whether an agent followed the right process, called the right tools, respected domain rules, and left the external system in a correct state.

For teams deploying agents in consequential settings, that question is not academic. A support agent that sounds polite but skips authentication, calls the wrong account tool, or leaves a ticket in an invalid state can cause real harm even when the final reply reads well. Building dependable agents therefore requires a workflow closer to software engineering than to prompt tuning alone: define the behaviour you expect, run the agent against those expectations, inspect what failed, fix the agent or policy, and repeat until the constraints hold. Without tooling that supports that loop, teams cannot move from ``it worked once in the playground'' to justified confidence that an agent will behave correctly under repeated, varied interaction.

Existing evaluation approaches only partly support that workflow. Traditional benchmarks and LLM-as-a-judge methods such as HELM, G-Eval, and MT-Bench are valuable for comparing models on isolated inputs and final outputs \cite{liang2022holistic,liu2023geval,zheng2023judging}, but they hide failures that live in tool traces and application state: a final answer can look reasonable while the agent skipped a required step or changed state incorrectly. Recent agent benchmarks such as AgentBench, WebArena, $\tau$-bench, and AppWorld have made trace- and state-aware evaluation more visible \cite{liu2023agentbench,zhou2023webarena,yao2024taubench,trivedi2024appworld}, yet they are fixed leaderboards rather than reusable infrastructure for a domain team writing its own behavioural contracts. Practical libraries and observability tools help developers run checks and inspect traces, but they still make it awkward to express a full multi-turn scenario (with output, tool, and state checks together) and to get precise diagnostics when something fails. Chapter~\ref{chap:related-work} analyses those limitations in detail.

This project addresses that need with \emph{agent-spec-kit}: a framework for scenario-driven evaluation of LLM agents. Its central bet is that reliability comes from \emph{declarative} behavioural specifications that a team can read, run, and iterate on. A scenario reads like a conversation script: user messages, checks on the agent's reply and tool use after each turn, and optional state oracles, all in one place. When an assertion fails, the framework reports which expectation broke and where in the trace it applied, so developers can fix the agent quickly rather than reconstruct the episode from a raw trace or aggregate score. The approach draws on software-testing ideas such as test oracles, repeated trials, and failure simplification \cite{claessen2000quickcheck,barr2015oracle,zeller2002simplifying,ribeiro2020checklist}. The goal is to make desired agent behaviour explicit, runnable, and inspectable, so failures are specified before they happen, discovered systematically, and preserved as regression tests for the next change.

\section{Contributions}

This project makes the following high-level contributions.

\begin{itemize}
    \item \textbf{A declarative, co-located scenario language for agent behaviour.} The project provides a way to define agent evaluations as executable scenarios rather than as isolated prompt--output pairs scattered across datasets, graders, and runner config. A scenario reads as a conversation script: user messages, checks on the agent's reply and tool use after each turn, optional environment steps, and repeat counts live in one place. That makes behavioural contracts easier for teammates to read, review, and change as product rules evolve.

    \item \textbf{Support for evaluating multiple surfaces of agent behaviour.} Instead of treating the final natural-language response as the only evaluation target, the framework supports checks over the user-facing output, intermediate tool interactions, and resulting application state. This reflects the observation from recent agent benchmarks that correctness in tool-using systems depends on trajectories and state changes, not only final answers \cite{yao2024taubench,trivedi2024appworld,lu2024toolsandbox}.

    \item \textbf{Assertion-local diagnostics and a testing-oriented failure workflow.} The project brings ideas from software testing into LLM agent evaluation, including repeated trials, generative variation of scenarios, simplification of failing cases, and extraction of regression tests. When an assertion fails, the framework reports which expectation broke, where in the trace it applied, and structured expected versus actual detail, so developers can fix the agent or policy without reconstructing the episode from a raw trace or aggregate score alone.

    \item \textbf{A practical evaluation workflow for domain-specific agents.} The framework is designed around the spec--run--inspect--fix loop that teams need when building agents with rules, tools, and stateful side effects. The contribution is not a new benchmark alone, but reusable infrastructure for constructing domain-specific evaluations where success criteria are explicit, rerunnable, and comparable across experiments.

    \item \textbf{An empirical evaluation of the framework on agent tasks.} The project evaluates the framework by applying it to realistic tool-using agent scenarios and analysing the kinds of failures that are found. The evaluation focuses not only on aggregate pass rates, but also on whether the framework helps expose distinct failure modes and turn them into reproducible regression cases.
\end{itemize}

Overall, this project contributes a developer-facing evaluation framework for LLM agents. Its central claim is that reliable agent development requires more than scalar benchmark scores or LLM-judged final answers: it requires declarative behavioural specs that a team can read, run repeatedly, and debug precisely when they fail.

